% sections/projets.tex
% Projets Scientifiques et Techniques — 3 projets maximum (contrainte 1 page)
% Chaque bullet utilise la formule XYZ ou STAR.
% Prioriser : pensée algorithmique, conception système, méthodologie scientifique.

\section{Projets Scientifiques et Techniques}

% --- Projet 1 : FLOW (matériel + IoT + prix national) ---
\cvitem{\textbf{FLOW} -- Dispositif IoT d'économie d'eau \newline \small{ESP32, C++, PCB, SolidWorks, Fusion 360} \newline \small{Fév. 2022 -- Déc. 2024}}{%
  \textbf{Contexte :} Conception d'un dispositif électromécanique et IoT pour réduire la consommation d'eau sous la douche.
  \begin{itemize}
    \item Architecturé et développé un système intégrant capteurs de débit, actionneurs de vanne, communication sans fil et PCB personnalisée, permettant une réduction de 60\% de la consommation d'eau.
    \item Lauréat, Heineken Green Challenge Mexique 2023 -- Top 20 des projets nationaux sélectionnés.
  \end{itemize}
}

% --- Projet 2 : ICPC (algorithmes, compétition internationale) ---
\cvitem{\textbf{ICPC} -- Concours International de Programmation \newline \small{C++, Algorithmes, Structures de données} \newline \small{Août 2024 -- Présent}}{%
  \textbf{Contexte :} Participation au concours universitaire ICPC, éditions 2024, 2025 et 2026.
  \begin{itemize}
    \item Classé 3\textsuperscript{e} au niveau Faculté UASLP, édition 2024 ; participant aux éditions 2025 et 2026.
    \item Développé et optimisé des solutions algorithmiques en C++ (graphes, programmation dynamique, géométrie computationnelle) dans un contexte compétitif chronométré.
  \end{itemize}
}

% --- Projet 3 : Hackathons 2024 (ML, produit, résultat en compétition) ---
\cvitem{\textbf{Hackathons 2024} -- NASA Space Apps \& Hackathon Magna \newline \small{Flutter, React, Python, ML} \newline \small{Oct. -- Nov. 2024}}{%
  \textbf{Contexte :} Deux hackathons successifs en 2024 avec des équipes pluridisciplinaires.
  \begin{itemize}
    \item NASA Space Apps Challenge Xalapa 2024 : Conçu une application de génération musicale à partir des données spectrographiques du télescope spatial James Webb.
    \item Hackathon Magna UASLP 2024 (Mention Honorable) : Développé une plateforme web/mobile cross-platform avec modèle de ML pour la prédiction de turnover employé (Flutter, React, Python).
  \end{itemize}
}

% TODO: 3 projets retenus sur 2 pages de matériel brut. Si contrainte de place, réduire les bullets à 1 par projet.
% TODO: Ajouter lien GitHub ou démo si disponible pour chaque projet.
